\documentclass[10pt,a4paper]{amsart}
\usepackage{amsmath,amssymb,amsthm,graphicx}
\usepackage{verbatim}
\newenvironment{reference}{\comment}{\endcomment}
%\newenvironment{reference}{}{}
\newenvironment{slogan}{\comment}{\endcomment}
\newenvironment{history}{\comment}{\endcomment}

% For commutative diagrams we use Xy-pic
\usepackage[all]{xy}

% We use 2cell for 2-commutative diagrams.
\xyoption{2cell}
% Theorem environments.
%
\theoremstyle{plain}
\newtheorem{theorem}[subsection]{Theorem}
\newtheorem{proposition}[subsection]{Proposition}
\newtheorem{lemma}[subsection]{Lemma}

\theoremstyle{definition}
\newtheorem{definition}[subsection]{Definition}
\newtheorem{example}[subsection]{Example}
\newtheorem{exercise}[subsection]{Exercise}
\newtheorem{situation}[subsection]{Situation}

\theoremstyle{remark}
\newtheorem{remark}[subsection]{Remark}
\newtheorem{remarks}[subsection]{Remarks}

\numberwithin{equation}{subsection}

% Macros
%
\def\lim{\mathop{\mathrm{lim}}\nolimits}
\def\colim{\mathop{\mathrm{colim}}\nolimits}
\def\Spec{\mathop{\mathrm{Spec}}}
\def\Hom{\mathop{\mathrm{Hom}}\nolimits}
\def\Ext{\mathop{\mathrm{Ext}}\nolimits}
\def\SheafHom{\mathop{\mathcal{H}\!\mathit{om}}\nolimits}
\def\SheafExt{\mathop{\mathcal{E}\!\mathit{xt}}\nolimits}
\def\Sch{\mathit{Sch}}
\def\Mor{\mathop{\mathrm{Mor}}\nolimits}
\def\Ob{\mathop{\mathrm{Ob}}\nolimits}
\def\Sh{\mathop{\mathit{Sh}}\nolimits}
\def\NL{\mathop{N\!L}\nolimits}
\def\CH{\mathop{\mathrm{CH}}\nolimits}
\def\proetale{{pro\text{-}\acute{e}tale}}
\def\etale{{\acute{e}tale}}
\def\QCoh{\mathit{QCoh}}
\def\Ker{\mathop{\mathrm{Ker}}}
\def\Im{\mathop{\mathrm{Im}}}
\def\Coker{\mathop{\mathrm{Coker}}}
\def\Coim{\mathop{\mathrm{Coim}}}

% Boxtimes
%
\DeclareMathSymbol{\boxtimes}{\mathbin}{AMSa}{"02}

%
% Macros for moduli stacks/spaces
%
\def\QCohstack{\mathcal{QC}\!\mathit{oh}}
\def\Cohstack{\mathcal{C}\!\mathit{oh}}
\def\Spacesstack{\mathcal{S}\!\mathit{paces}}
\def\Quotfunctor{\mathrm{Quot}}
\def\Hilbfunctor{\mathrm{Hilb}}
\def\Curvesstack{\mathcal{C}\!\mathit{urves}}
\def\Polarizedstack{\mathcal{P}\!\mathit{olarized}}
\def\Complexesstack{\mathcal{C}\!\mathit{omplexes}}
% \Pic is the operator that assigns to X its picard group, usage \Pic(X)
% \Picardstack_{X/B} denotes the Picard stack of X over B
% \Picardfunctor_{X/B} denotes the Picard functor of X over B
\def\Pic{\mathop{\mathrm{Pic}}\nolimits}
\def\Picardstack{\mathcal{P}\!\mathit{ic}}
\def\Picardfunctor{\mathrm{Pic}}
\def\Deformationcategory{\mathcal{D}\!\mathit{ef}}
\usepackage{cleveref}
\crefname{originallemma}{original lemma}{original lemmas}
\Crefname{originallemma}{Original lemma}{Original lemmas}
\makeatletter
\newlabel{lemma-module-principal-parts}{{09CK}{source}{}{sourcepointer-1}{}}
\newlabel{lemma-module-principal-parts@cref}{{[originallemma][0][]09CK}{[1][0][]source}{}{}{}}
\makeatother
\title{Principal parts of order k as a quotient of the base change by the power of the diagonal ideal}
\author{The Stacks Project Authors}
\date{stacks-project:algebra.tex:local-snapshot-20261012; GNU FDL 1.2 or later}
\begin{document}
\maketitle
\noindent\textit{Editorial note.} All delivered mathematical text is the author's own, copied from the pinned Stacks Project chapter file algebra.tex; no statement or proof step was rewritten, reordered or removed.\par\medskip
\noindent\textit{Editorial note.} Preamble provenance, as required for this source: the frozen snapshot stages the chapter's own preamble file beside it (preamble.tex, 217 lines, 8116 bytes, the same shared Stacks-wide preamble the published chow and etale-cohomology cards use); the first line of the pinned chapter file is \textbackslash{}input\{preamble\}, so no preamble text lives inside the chapter itself. The card takes that file's lines 8-18 (the verbatim package and the three \textbackslash{}comment-based environments reference/slogan/history at lines 9-12, section 13.2, together with the Xy-pic setup) and its lines 154 to the runtime end of the file (the native theorem declarations and mathematical macros). Not one character was changed, no package was added, no macro was defined, and the excerpt serves only as preamble, never as a span coordinate.\par\medskip
\noindent\textit{Editorial note.} The same-chapter pointers are located at their \textbackslash{}label lines inside the pinned chapter file; the cross-chapter pointers are located in the project's own tag list staged in the snapshot (sources/stacks-chow/tags.txt), at the verbatim lines carrying their canonical tags.\par\medskip
\noindent\textit{Editorial note.} The counted proof ends by leaving the verification of the inverse identity to the reader; that sentence is kept verbatim inside the span and both maps are written out.\par\medskip
\noindent\textit{Editorial note.} The counted proof refers to a same-chapter statement by label; it is recorded as a navigation stub pointed at its \textbackslash{}label line.\par\medskip
\noindent\textit{Allowed standard prerequisites (editorial).}\begin{itemize}
\item The universal differential operator of order k. Used throughout the counted proof (lines 35160-35184). Supplies the map that the isomorphism has to identify with the class of one tensor m.
\item The diagonal ideal as the kernel of the multiplication map. Used in the statement and in the counted proof. Provides the ideal whose (k+1)-st power cuts out the target of the isomorphism.
\end{itemize}\medskip
\ifdefined\hypertarget\hypertarget{sourcepointer-1}{}\fi
\noindent\textit{Source reference.} Original lemma 09CK: sources/stacks-algebra/pinned-original-source.tex, source lines 34918--34918. Navigation pointer for \textbackslash{}ref\{lemma-module-principal-parts\} used inside the counted span (line 35172); the argument that consumes it is retained in the span.\par\medskip
\begin{lemma}
\label{lemma-principal-parts-diagonal}
Let $R \to S$ be a ring map. Let $M$ be an $S$-module.
Let $J = \Ker(S \otimes_R S \to S)$ be the kernel of the multiplication map. 
There is a canonical isomorphism of $S$-modules
$$
P^k_{S/R}(M) \longrightarrow (S \otimes_R M)/J^{k + 1}(S \otimes_R M)
$$
where $s \in S$ acts on the target via multiplication by $s \otimes 1$
and such that the universal differential operators of order $k$
to the map given by $m \mapsto \text{class of }1 \otimes m$.
\end{lemma}

\begin{proof}
Consider the map $T : M \to S \otimes M$, $m \mapsto 1 \otimes m$
Since $T$ is $R$-linear and since
\begin{align*}
g_0g_1\ldots g_k T(m)
- \sum g_0 \ldots \hat g_i \ldots g_k T(g_im) + \ldots
+(-1)^{k + 1}T(g_0\ldots g_km) \\
=
\prod_{i = 0,\ldots,k} (g_i \otimes 1 - 1 \otimes g_i) 1 \otimes m
\end{align*}
for $m \in M$ and $g_0, \ldots, g_k \in S$, we conclude that
the rule $m \mapsto \text{class of }1 \otimes m$ in the statement
of the lemma is a differential operator of order $k$
(see proof of Lemma \ref{lemma-module-principal-parts}).
By the universal property of $P^k_{S/R}(M)$ we obtain the
arrow in the statement of the lemma.
On the other hand, if $D : M \to N$ is a differential operator
of order $k$, then we can consider the $S$-linear map
$L : S \otimes_R M \to N$, $g \otimes m \mapsto gD(m)$.
The reader checks, by a computation similar to the one above
and using that $J$ is generated by elements of the form
$g \otimes 1 - 1 \otimes g$, that
$L$ annihilates $J^{k + 1}(S \otimes_R M)$.
In particular, if we apply this to the universal differential
operator $D_{univ} : M \to P^k_{S/R}(M)$, then we obtain an
$S$-linear map $(S \otimes_R M)/J^{k + 1}(S \otimes_R M) \to P^k_{S/R}(M)$.
We leave it to the reader to see that this map is the inverse to
the map in the statement of the lemma.
\end{proof}

\end{document}
